\chapter{Үр дүн ба үнэлгээ}

Энэ бүлэгт системийн чанарыг зөвхөн өндөр тоон үзүүлэлтээр сурталчлах бус, тухайн үзүүлэлтийг ямар нөхцөлд авсан, ямар хэмжээнд тайлбарлах ёстой, мөн ямар хязгаарлалтыг дагалдуулан унших хэрэгтэйг хамтад нь авч үзэв. Ийм байр суурь нь академик шударга тайлагналын үндэс болно.

\section{Туршилтын орчин}

Бид ML загварыг сургах болон системийг туршихдаа дараах орчинд ажиллуулсан:

\begin{table}[H]
\centering
\caption{Туршилтын орчны тохиргоо}
\label{tab:test_env}
\begin{longtable}{|l|l|l|}
\hline
\textbf{Бүрэлдэхүүн} & \textbf{Сургалт} & \textbf{Inference} \\
\hline
CPU & AMD Ryzen 7 5800X & Intel i5-12400 \\
\hline
GPU & NVIDIA RTX 3080 & NVIDIA RTX 3060 \\
\hline
RAM & 32GB DDR4 & 16GB DDR4 \\
\hline
OS & Ubuntu 22.04 & Ubuntu 22.04 \\
\hline
Python & 3.10.12 & 3.10.12 \\
\hline
PyTorch & 2.0.1+cu118 & 2.0.1+cu118 \\
\hline
\end{longtable}
\end{table}

Серверийн туршилтыг Docker хэрэгслээр контейнержуулж, MongoDB-тэй холбогдох асуудлыг урьдчилан шийдвэрлэсний үндсэн дээр тогтвортой ажиллах орчныг бүрдүүлсэн.

\section{Загварын нарийвчлалын хэмжүүрүүд}

Объект илрүүлэлтийн загварын үнэлгээнд түгээмэл ашиглагддаг нийтлэг параметрүүдыг тус тус хэмжсэн.

\subsection{Үнэлгээний хэмжүүрүүдийн тодорхойлолт}

\begin{itemize}
    \item \textbf{Precision (Нарийвчлал):} Нийт эерэг гэж таамагласан объектуудаас хэд нь үнэхээр тухайн ангиллын бараа байсныг илтгэнэ.
    \begin{equation}
        Precision = \frac{TP}{TP + FP}
    \end{equation}
    
    \item \textbf{Recall (Мэдрэг байдал):} Бодит зураг дээр байгаа нийт объектуудын хэдэн хувийг нь зөв илрүүлж чадсаныг заана.
    \begin{equation}
        Recall = \frac{TP}{TP + FN}
    \end{equation}
    
    \item \textbf{F1 Score:} Precision болон Recall-ийн гармоник дундаж.
    \begin{equation}
        F1 = 2 \times \frac{Precision \times Recall}{Precision + Recall}
    \end{equation}
    
    \item \textbf{mAP (Mean Average Precision):} Бүх ангиллын Average Precision-ий дундаж.
    \begin{equation}
        mAP = \frac{1}{N} \sum_{i=1}^{N} AP_i
    \end{equation}
    
    \item \textbf{IoU (Intersection over Union):} Таамагласан болон бодит bounding box-ийн давхцлын хэмжүүр.
    \begin{equation}
        IoU = \frac{Area_{intersection}}{Area_{union}}
    \end{equation}
\end{itemize}

\subsection{Ерөнхий үр дүн}

\begin{table}[H]
\centering
\caption{Загварын ерөнхий үнэлгээний үр дүн}
\label{tab:overall_results}
\begin{longtable}{|l|c|}
\hline
\textbf{Хэмжүүр} & \textbf{Утга} \\
\hline
mAP@0.5 (mAP50) & 92.3\% \\
\hline
mAP@0.5:0.95 (mAP50-95) & 78.6\% \\
\hline
Precision & 94.1\% \\
\hline
Recall & 89.7\% \\
\hline
F1 Score & 91.8\% \\
\hline
\end{longtable}
\end{table}

\subsection{Үр дүнгийн найдвартай байдлын тайлбар}

Энэ бүлэгт тайлагнасан \texttt{mAP50 = 92.3\%}, \texttt{Precision = 94.1\%}, \texttt{Recall = 89.7\%} үзүүлэлтүүд нь энэхүү ажлын хүрээнд хадгалсан нэг үндсэн сургалтын туршилтын эцсийн үнэлгээ болно. Олон random seed ашигласан 3--5 давталттай сургалтын дундаж, стандарт хазайлт тооцсон тусгай туршилт энэ хувилбарын ажлын хүрээнд хийгдээгүй тул статистикийн утгаар ``\(92.3 \pm x\)\%'' хэлбэрийн дүгнэлт хийх боломжгүй. Иймээс эдгээр үр дүнг системийн боломжийг харуулсан \textit{эмпирик гол үзүүлэлт} гэж үзэх нь зохистой бөгөөд давтагдах чанарыг баталгаажуулах олон давталттай туршилтыг цаашдын ажлын чухал чиглэл гэж үзэв.

Гэсэн хэдий ч тухайн үзүүлэлтүүдийг нэг ижил тестийн багц, нэг ижил үнэлгээний pipeline, нэг ижил inference тохиргоогоор тооцсон тул хоорондын дотоод харьцуулалт хүчинтэй байна. Өөрөөр хэлбэл, ангилал хоорондын ялгаа, error pattern, audit decision-д үзүүлж буй нөлөөллийг энэхүү туршилтын хүрээнд шинжлэх нь үндэслэлтэй юм.

\subsection{Ангилал тус бүрийн үр дүн}

\begin{table}[H]
\centering
\caption{Ангилал тус бүрийн нарийвчлалын үнэлгээ}
\label{tab:class_results}
\begin{longtable}{|l|c|c|c|c|}
\hline
\textbf{Ангилал} & \textbf{Precision} & \textbf{Recall} & \textbf{F1} & \textbf{AP@0.5} \\
\hline
Coca-Cola 500ml & 96.2\% & 93.5\% & 94.8\% & 95.1\% \\
\hline
Pepsi 500ml & 94.8\% & 91.2\% & 93.0\% & 93.4\% \\
\hline
Fanta 500ml & 95.1\% & 90.8\% & 92.9\% & 93.2\% \\
\hline
Sprite 500ml & 93.7\% & 89.4\% & 91.5\% & 91.8\% \\
\hline
Bon Aqua 500ml & 94.5\% & 90.1\% & 92.2\% & 92.6\% \\
\hline
Vittel 500ml & 91.8\% & 86.3\% & 88.9\% & 89.2\% \\
\hline
Rich Orange 1L & 93.2\% & 88.7\% & 90.9\% & 91.1\% \\
\hline
Dobry Apple 1L & 92.9\% & 87.4\% & 90.1\% & 90.4\% \\
\hline
\textbf{Дундаж} & \textbf{94.1\%} & \textbf{89.7\%} & \textbf{91.8\%} & \textbf{92.3\%} \\
\hline
\end{longtable}
\end{table}

Confusion Matrix-ын дүнгээс харахад хамгийн тогтвортой ангиллууд нь \texttt{Coca-Cola}, \texttt{Bon Aqua} зэрэг өнгө, логотип, савлагааны silhouette нь ялгаралт өндөр бүтээгдэхүүнүүд байв. Харин \texttt{Vittel 500ml}, \texttt{Dobry Apple 1L}, \texttt{Rich Orange 1L} зэрэг ангиллуудын нарийвчлал харьцангуй бага гарсан нь дараах бодит шалтгаантай уялдаж байна. Нэгдүгээрт, тунгалаг эсвэл цайвар шошготой ус, ундааны савнууд shelf дээр зэрэгцэн байрлах үед савлагааны ерөнхий хэлбэр ойролцоо тул ангиллын ялгаа багасдаг. Хоёрдугаарт, 1 литрийн тетра савлагаатай жүүсүүдэд брэндийн үндсэн layout төстэй байсан нь жижиг resolution дээр логотипийн ялгааг багасгасан. Гуравдугаарт, хэсэгчлэн халхлагдсан болон glare үүссэн зурагт загвар дүрсний texture-ээс илүү өнгөний ерөнхий тархалтад түшиглэн шийдвэр гаргах хандлага ажиглагдсан.

Практик талаас нь авч үзвэл, энэхүү андуурал нь зөвхөн detector-ийн түвшний ангиллын алдаа биш, харин аудитын дараагийн шатанд ``аль SKU дутуу байна вэ'' гэсэн шийдвэрт нөлөөлөх эрсдэлтэй. Иймээс confusion matrix-ийг дан ганц model quality-ийн үзүүлэлт бус, бизнесийн шийдвэрт дам нөлөө үзүүлэх эрсдэлийн зураглал гэж үзэж тайлбарлах шаардлагатай.

\subsection{Precision-Recall муруй}

Precision-Recall үнэлгээгээр бүх ангиллын дундаж муруй тогтвортой буурч, өндөр талбай эзэлсэн бөгөөд энэ нь \texttt{mAP@0.5 = 92.3\%} үзүүлэлттэй нийцэж байна.

\subsection{Confusion Matrix}

Confusion Matrix нь загварын ангилал тус бүрийн гүйцэтгэлийг харуулдаг. Диагональ дээрх тоонууд нь зөв илрүүлэлт, диагоналиас гадуурх тоонууд нь андуурсан илрүүлэлтийг илэрхийлнэ.

\begin{table}[H]
\centering
\caption{Confusion Matrix-ын товч үзүүлэлт}
\label{tab:confusion_summary}
\begin{longtable}{|l|c|c|c|}
\hline
\textbf{Ангилал} & \textbf{Зөв} & \textbf{Андуурсан} & \textbf{Нарийвчлал} \\
\hline
Coca-Cola & 412 & 8 & 98.1\% \\
\hline
Pepsi & 389 & 12 & 97.0\% \\
\hline
Fanta & 356 & 15 & 96.0\% \\
\hline
Sprite & 341 & 18 & 95.0\% \\
\hline
Bon Aqua & 398 & 10 & 97.5\% \\
\hline
Vittel & 287 & 22 & 92.9\% \\
\hline
Rich Orange & 334 & 16 & 95.4\% \\
\hline
Dobry Apple & 330 & 18 & 94.8\% \\
\hline
\end{longtable}
\end{table}

\section{Системийн хурдны үнэлгээ}

Зөвхөн нарийвчлал сайн байхаас гадна, боловсруулалт хурдан байх нь системд нэвтрүүлэх том шалгуур байдаг.

\subsection{Боловсруулах хугацааны хэмжилт}

\begin{table}[H]
\centering
\caption{Системийн боловсруулах хурдны үнэлгээ}
\label{tab:speed_results}
\begin{longtable}{|l|c|c|c|}
\hline
\textbf{Үе шат} & \textbf{Дундаж (ms)} & \textbf{Хамгийн бага} & \textbf{Хамгийн их} \\
\hline
Зураг урьдчилан боловсруулах & 15 & 12 & 23 \\
\hline
YOLOv8 inference & 185 & 156 & 234 \\
\hline
NMS (Non-Max Suppression) & 12 & 8 & 18 \\
\hline
Үр дүн боловсруулах & 8 & 5 & 15 \\
\hline
MongoDB хадгалах & 25 & 18 & 45 \\
\hline
\textbf{Нийт} & \textbf{245} & \textbf{199} & \textbf{335} \\
\hline
\end{longtable}
\end{table}

Хугацааны задаргаанаас хамгийн их ачаалал \texttt{YOLOv8 inference} шатанд (185ms) төвлөрч, бусад шатны нийлбэр хугацаа харьцангуй бага (60ms) байв.

\subsection{FPS (Frames Per Second) хэмжилт}

\begin{table}[H]
\centering
\caption{FPS хэмжилт янз бүрийн төхөөрөмж дээр}
\label{tab:fps_results}
\begin{longtable}{|l|c|c|}
\hline
\textbf{Төхөөрөмж} & \textbf{FPS} & \textbf{Latency (ms)} \\
\hline
NVIDIA RTX 3080 & 28.5 & 35 \\
\hline
NVIDIA RTX 3060 & 18.2 & 55 \\
\hline
NVIDIA GTX 1660 & 12.4 & 81 \\
\hline
CPU only (i5-12400) & 3.2 & 312 \\
\hline
\end{longtable}
\end{table}

Иймд энэхүү хугацаа нь хүлээгдэх хугацааны шаардлагыг бүрэн хангаж байгаа боловч энд тайлагнасан \texttt{245ms} нь backend inference pipeline-ийн дундаж хугацаа гэдгийг ялгаж ойлгох хэрэгтэй. Хэрэглэгчийн бодит туршлагад үүн дээр зураг upload, мобайл сүлжээний latency, серверийн ачаалал, response rendering зэрэг нэмэлт хүчин зүйлүүд нөлөөлнө. Тиймээс ``245ms''-ийг хэрэглэгч яг ийм хугацаанд хариу авна гэсэн шууд UI latency биш, харин сервер талын bottleneck хүлээн зөвшөөрөхүйц түвшинд байгааг харуулсан инженерийн үзүүлэлт гэж тайлбарлах нь зөв.

\section{Аудитын алгоритмын үнэлгээ}

Аудитын хөдөлгүүрийн нарийвчлалыг үнэлэхийн тулд 50 удаагийн туршилтыг янз бүрийн нөхцөлд хийсэн.

\begin{table}[H]
\centering
\caption{Аудитын үр дүнгийн задаргаа}
\label{tab:audit_results}
\begin{longtable}{|l|c|c|c|}
\hline
\textbf{Статус} & \textbf{Зөв} & \textbf{Буруу} & \textbf{Нарийвчлал} \\
\hline
PASS & 18 & 1 & 94.7\% \\
\hline
WARNING & 15 & 2 & 88.2\% \\
\hline
FAIL & 12 & 2 & 85.7\% \\
\hline
\textbf{Нийт} & \textbf{45} & \textbf{5} & \textbf{90.0\%} \\
\hline
\end{longtable}
\end{table}

Энд нийт аудитын зөв ангилалтын хувь \(90.0\%\) байгаа боловч статус тус бүр ижил түвшинд ажиллаагүй нь анхаарал татна. Ялангуяа \texttt{FAIL} статусын нарийвчлал \texttt{PASS}-аас бага (\(85.7\%\) vs. \(94.7\%\)) байгаа нь аудитын алгоритмын decision rule-ээс илүү detector-ийн алдаа хэрхэн дамжин нөлөөлж байгааг харуулж байна.

\texttt{FAIL} шийдвэр нь ихэвчлэн бараа дутсан эсвэл тохирлын хувь босгоос доош орсон нөхцөлд гардаг. Ийм үед detector-ийн \texttt{False Negative} алдаа, өөрөөр хэлбэл shelf дээр байгаа барааг илрүүлж чадаагүй тохиолдол, аудитын тооцоонд шууд ``дутуу бараа'' болж бүртгэгдэнэ. Үүний улмаас бодит байдал дээр \texttt{WARNING} эсвэл бүр \texttt{PASS} байх нөхцөлийг \texttt{FAIL} болгон хатуу ангилах эрсдэл нэмэгддэг. Харин \texttt{PASS} статус нь илүү ``өршөөлтэй'' бүтэцтэй: detector нэгж түвшний бага алдаа гаргасан ч нийт тохирлын хувь босгоос дээш үлдэх магадлал өндөр байдаг.

Иймээс аудитын хөдөлгүүрийн гүйцэтгэлийг дангаар нь бус, \textit{detector + decision logic} гэсэн цуваа системийн гүйцэтгэл гэж авч үзэх нь зөв. Цаашид detector-ийн confidence score-ыг decision layer-д оруулж, босго орчим үр дүнг ``баталгаажуулах шаардлагатай'' гэсэн завсрын ангилалд шилжүүлэх нь \texttt{FAIL} ангиллын буруу сөрөг шийдвэрийг бууруулах боломжтой.

\section{Алдааны шинжилгээ (Error Analysis)}

\subsection{Түгээмэл алдааны төрлүүд}

Туршилтын явцад дараах алдаанууд ажиглагдсан:

\begin{enumerate}
    \item \textbf{False Positive (Буруу эерэг):} Байхгүй барааг байна гэж таньсан тохиолдол --- 4.2\%
    \item \textbf{False Negative (Буруу сөрөг):} Байгаа барааг илрүүлж чадаагүй тохиолдол --- 6.1\%
    \item \textbf{Төстэй бараа андуурах:} Coca-Cola болон Pepsi-г андуурах --- 2.3\%
    \item \textbf{Хагас халхлагдсан бараа:} 50\%-иас дээш халхлагдсан үед --- 8.5\%
\end{enumerate}

Алдааны тархалтаас харахад хамгийн том хувь нь \texttt{False Negative} болон \texttt{халхлагдсан бараа}-тай холбоотой байсан тул цаашид өгөгдөл баяжуулалт болон camera angle стандартчилал хийх шаардлагатай. Энэ нь detector-ийн алдаа зөвхөн нэгж илрүүлэлтийн түвшинд үлдэхгүй, харин аудитын дүгнэлтэд ``дутуу бараа'' гэсэн буруу дохио үүсгэх эрсдэлтэйг харуулна. Иймээс error analysis-ийг model debugging төдийгөөр ойлгох бус, decision reliability-ийн үндсэн эх сурвалж гэж үзэв.

\subsection{Алдаа бууруулах боломжит арга замууд}

\begin{longtable}{|p{3.2cm}|p{3.8cm}|p{5.2cm}|}
\caption{Түгээмэл алдаанд тохирох техникийн сайжруулалтын чиглэл}
\label{tab:error_mitigation} \\
\hline
\textbf{Алдааны төрөл} & \textbf{Ажиглагдсан шалтгаан} & \textbf{Сайжруулалтын боломж} \\
\hline
Хагас халхлалт & Нэг барааны урд өөр бараа давхарласан, shelf нягт өрөлттэй байсан & Олон өнцгөөс зураг авах протокол нэвтрүүлэх, partial occlusion augmentation нэмэх, planogram-тэй хослуулсан prior ашиглах \\
\hline
Бүдэг гэрэл & Contrast буурч, шошгоны texture ялгарахгүй болсон & Brightness/contrast augmentation өргөтгөх, зураг авах minimum lux заавар гаргах, auto-exposure шалгалттай capture workflow хэрэглэх \\
\hline
Хэт гэрэл, glare & Савлагаан дээр specular reflection үүсч логотип халхлагдсан & Glare-rich samples-ийг датасетэд нэмэх, preprocessing-д highlight suppression турших, retake guidance өгөх \\
\hline
Төстэй савлагаа & Өнгө, silhouette ойролцоо SKU-ууд андуурагдсан & Fine-grained class balancing хийх, илүү өндөр нягтралтай crop-based re-classification layer нэмэх \\
\hline
False Negative & Жижиг объект, хэт нягт өрөлт, халхлалт & Multi-scale inference турших, hard example mining хийх, low-confidence detections-ийг audit layer-д болгоомжтой жинлэх \\
\hline
\end{longtable}

Эдгээр арга замын заримаас энэхүү судалгаанд зөвхөн өгөгдөл баяжуулалт болон зураг авах стандартчиллын түвшинд хязгаарлагдмал хэлбэрээр тусгасан. Харин planogram-aware inference, multi-view capture, хоёр шаттай detector-classifier pipeline зэрэг нь цаашдын судалгааны илүү гүн чиглэлд хамаарна.

\subsection{Гэрэлтүүлгийн нөлөө}

\begin{table}[H]
\centering
\caption{Гэрэлтүүлгийн нөхцөл дэх нарийвчлал}
\label{tab:lighting_results}
\begin{longtable}{|l|c|c|}
\hline
\textbf{Гэрэлтүүлгийн нөхцөл} & \textbf{mAP50} & \textbf{Recall} \\
\hline
Хэвийн гэрэлтүүлэг & 94.5\% & 92.1\% \\
\hline
Бүдэг гэрэл & 88.3\% & 84.6\% \\
\hline
Хэт гэрэлтэй & 90.1\% & 86.8\% \\
\hline
Холимог гэрэл & 91.2\% & 88.4\% \\
\hline
\end{longtable}
\end{table}

Гэрэлтүүлгийн туршилтын үр дүнгээс харахад бүдэг гэрэлтэй нөхцөлд \texttt{mAP50} болон \texttt{Recall} хоёулаа хамгийн их буурсан байна. Энэ нь shelf зураг дахь ялгаралт багассанаар contour, label texture, өнгөний хил зааг бүдгэрч байгаатай холбоотой. Хэт гэрэлтэй нөхцөлд гарсан бууралт нь бүдэг гэрлийнхээс бага байгаа боловч ус, шингэн савлагааны гадаргуу дээр үүсэх glare нь detector-ийн box localization болон class confidence-д сөргөөр нөлөөлсөн.

Практик түвшинд авч үзвэл, энэхүү үр дүн нь Random Brightness зэрэг augmentation нь гэрэлтүүлгийн өөрчлөлтөд тодорхой хэмжээнд тэсвэртэй болгосон ч бодит дэлгүүрийн орчин дахь бүх lighting condition-ийг бүрэн хамарч чадаагүйг харуулж байна. Иймээс цаашид гэрэлтүүлгийн нөхцөлийг тусгайлан стратифай хийсэн dataset өргөтгөл, capture-time quality check, шаардлагатай тохиолдолд flash/retake guidance нэвтрүүлэх нь үр дүнтэй гэж үзэв.

\section{Бусад системүүдтэй харьцуулалт}

Zhang нар 2020 онд ``Automatic Shelf Monitoring and Product Recognition Using Deep Learning'' судалгаандаа дэлгүүрийн лангууны автомат хяналтын системийг хөгжүүлж, YOLOv5 ашиглан 89.2\% mAP нарийвчлал хүрсэн \cite{shelf_monitoring}. Манай системийн гол үзүүлэлт үүнээс өндөр гарсан боловч энэ нь ижил датасет, ижил annotation protocol, ижил hardware, ижил training schedule дээр хийсэн шууд benchmark биш гэдгийг тодотгох шаардлагатай. Иймээс 3.1\%-ийн ялгааг \textit{шууд superiority claim} биш, харин тухайн ажлуудын тайлагнасан хүрээнд харьцуулсан лавлагаа үзүүлэлт гэж үзэв.

Tonioni нар 2017 онд ``Planogram Compliance'' судалгаандаа лангууны барааны байршлыг шалгах системийг graph matching аргаар хэрэгжүүлсэн \cite{planogram_compliance}. Бидний систем одоогоор зөвхөн барааны тоо хэмжээг шалгадаг боловч, цаашид байршлын шалгалт нэмэх боломжтой.

\begin{table}[H]
\centering
\caption{Бусад объект илрүүлэлтийн загваруудтай харьцуулалт}
\label{tab:model_comparison}
\begin{longtable}{|l|c|c|c|c|}
\hline
\textbf{Загвар} & \textbf{mAP50} & \textbf{FPS} & \textbf{Хэмжээ (MB)} & \textbf{Параметр} \\
\hline
YOLOv5n & 89.1\% & 24.3 & 3.9 & 1.9M \\
\hline
YOLOv5s & 91.4\% & 18.7 & 14.1 & 7.2M \\
\hline
\textbf{YOLOv8n (Бидний)} & \textbf{92.3\%} & \textbf{28.5} & \textbf{6.3} & \textbf{3.2M} \\
\hline
YOLOv8s & 93.8\% & 21.2 & 22.5 & 11.2M \\
\hline
Faster R-CNN & 94.2\% & 5.1 & 108.5 & 41.1M \\
\hline
SSD MobileNet & 85.6\% & 35.2 & 19.3 & 6.8M \\
\hline
\end{longtable}
\end{table}

Хүснэгт \ref{tab:model_comparison}-д өгөгдсөн зарим тоо нь ижил датасет дээр дахин сургаж авсан дүн биш, харин тухайн загварын нийтлэгдсэн эсвэл түгээмэл тайлагнагддаг үзүүлэлттэй харьцуулсан тойм шинжтэй утгууд болно. Тиймээс энэхүү харьцуулалтын зорилго нь \texttt{YOLOv8n}-ийн \textit{capacity-speed trade-off}-ийг тайлбарлахад оршихоос, статистикийн хувьд шударга head-to-head benchmark батлахад оршихгүй. Энэхүү болгоомжлолыг харгалзан үзэхэд, манай орчинд \texttt{YOLOv8n} нь нарийвчлал болон хурдны хооронд практик хэрэгжилтэд тохирсон тэнцвэр үзүүлсэн гэж дүгнэв.

\section{Хэрэглэгчийн туршилт (User Testing)}

Системийг бодит хэрэглэгчдээр туршуулж, санал асуулга авсан. Туршилтад нийт \(n=15\) оролцогч хамрагдсан бөгөөд тэдгээрийг тухайн системийн боломжит хэрэглээтэй уялдуулан сонгосон.

\begin{longtable}{|p{3.5cm}|c|p{6cm}|}
\caption{Хэрэглэгчийн туршилтын оролцогчдын бүтэц}
\label{tab:user_test_participants} \\
\hline
\textbf{Бүлэг} & \textbf{Тоо} & \textbf{Тайлбар} \\
\hline
Аудитор / field staff & 6 & Зураг авч, shelf audit хийх үндсэн хэрэглэгчид \\
\hline
Дэлгүүрийн менежер & 4 & Тайлан, статистик, зөрүүний мэдээлэл ашиглах хэрэглэгчид \\
\hline
Судалгааны туслах / оюутан & 5 & UI ойлгомжтой байдал, onboarding шалгах зорилготой \\
\hline
\end{longtable}

Туршилтыг 2026 оны 2-р сард хоёр үе шаттай зохион байгуулсан. Нэг оролцогчид дунджаар 20--25 минутын хугацаанд дараах даалгавруудыг өгсөн. Туршилтын протоколын хүрээнд оролцогч бүрт богино зааварчилгаа өгч, дараа нь бие даан даалгавар гүйцэтгүүлэн, төгсгөлд нь Likert-scale санал асуулга авч qualitative ажиглалтыг тэмдэглэсэн.

\begin{longtable}{|p{0.9cm}|p{10.4cm}|}
\caption{Хэрэглэгчийн туршилтын даалгаврууд}
\label{tab:user_test_tasks} \\
\hline
\textbf{№} & \textbf{Даалгавар} \\
\hline
T1 & Системд нэвтэрч өөрт хуваарилагдсан кампанит ажлыг сонгох \\
\hline
T2 & Дэлгүүр, лангуу сонгож зураг авч сервер рүү илгээх \\
\hline
T3 & Илрүүлэлтийн үр дүнгээс PASS/WARNING/FAIL статусыг тайлбарлах \\
\hline
T4 & Өмнөх аудитын түүхээс нэг тайланг нээж шалгах \\
\hline
T5 & Менежерийн дүрээр статистик хуудас болон зөрүүний тайланг харах \\
\hline
\end{longtable}

\begin{table}[H]
\centering
\caption{Хэрэглэгчийн туршилтын үр дүн (n=15)}
\label{tab:user_testing}
\begin{longtable}{|l|c|}
\hline
\textbf{Асуулт} & \textbf{Дундаж оноо (1--5)} \\
\hline
Системийг ашиглахад хялбар байсан уу? & 4.3 \\
\hline
Илрүүлэлтийн хурд хангалттай байсан уу? & 4.6 \\
\hline
Үр дүнгийн нарийвчлал хангалттай байсан уу? & 4.1 \\
\hline
Аудитын статус ойлгомжтой байсан уу? & 4.5 \\
\hline
Систем бодит ажилд хэрэглэх боломжтой юу? & 4.2 \\
\hline
\textbf{Ерөнхий үнэлгээ} & \textbf{4.34} \\
\hline
\end{longtable}
\end{table}

Likert-scale асуулгын нийлбэр үнэлгээг System Usability Scale (SUS)-тай дүйцүүлэн хөрвүүлэхэд систем 81.7 оноо авсан. Brooke-ийн ангиллаар 80-аас дээш SUS нь ``excellent'' ангилалд багтдаг тул энэхүү системийн хэрэглэхэд хялбар байдал зохих түвшинд байна гэж үзэв. Гэхдээ энд ашигласан үнэлгээ нь \(n=15\) хэмжээтэй, зорилтот хэрэглэгч болон сургалтын туслахуудыг хамарсан convenience sample байсан тул бүх төрлийн retail хэрэглэгч рүү статистикийн утгаар ерөнхийшүүлэхэд болгоомжтой хандах шаардлагатай.

Qualitative ажиглалтаас харахад оролцогчдын дунд хоёр нийтлэг хэв шинж давтагдсан. Нэгдүгээрт, зураг авах болон үр дүн харах урсгалыг ``цөөн алхамтай, ойлгомжтой'' гэж үнэлэх хандлага давамгай байв. Нэг оролцогчийн тэмдэглэлд ``зураг аваад шууд статус гарч ирж байгаа нь ажлын урсгалыг таслахгүй байсан'' гэж дурдсан нь field workflow-той нийцэж буйг илтгэнэ. Хоёрдугаарт, зарим оролцогч \texttt{WARNING} ба \texttt{FAIL} статустай холбоотой тайлбар, ялангуяа зөрүүний шалтгааныг илүү дэлгэрэнгүй харах хэрэгцээ байгааг тэмдэглэсэн. Үүнийг ``анхааруулга яагаад гарсныг арай дэлгэрэнгүй хармаар байна'' гэсэн санал давтан нотолсон. Энэхүү qualitative дүгнэлт нь UI/UX сайжруулалтын дараагийн шатанд status explanation болон discrepancy visualization-ийг өргөтгөх шаардлагатайг илтгэнэ.

\section{Үр дүнгийн хүчин төгөлдөр байдлын тухай тайлбар}

Энэхүү үнэлгээний бүлгийн үр дүнг уншихдаа дотоод, гадаад, хэмжилтийн хүчин төгөлдөр байдлыг тусад нь авч үзэх шаардлагатай.

\begin{longtable}{|p{3.2cm}|p{3.9cm}|p{4.9cm}|}
\caption{Үнэлгээний хүчин төгөлдөр байдлын гол эрсдэлүүд}
\label{tab:validity_threats} \\
\hline
\textbf{Төрөл} & \textbf{Эрсдэл} & \textbf{Энэ ажилд хэрхэн тайлбарласан} \\
\hline
\endfirsthead
\hline
\textbf{Төрөл} & \textbf{Эрсдэл} & \textbf{Энэ ажилд хэрхэн тайлбарласан} \\
\hline
\endhead
\hline
\endfoot
Internal validity & Detector болон audit engine-ийн алдаа хоорондоо холилдох & 6.4-т audit шийдвэрийг detector-оос тусгаарлахгүйгээр цуваа систем гэж тайлбарласан \\
\hline
External validity & Датасет ба хэрэглэгчийн туршилт хязгаарлагдмал орчинд хийгдсэн & Улаанбаатарын хязгаарлагдмал дэлгүүр, \(n=15\) хэрэглэгчийн туршилт тул бүх retail орчинд шууд ерөнхийшүүлэхээс зайлсхийсэн \\
\hline
Construct validity & ``Хэрэглэхэд хялбар'' гэсэн ойлголтыг шууд нэг тоогоор хэмжих хязгаарлалтай & Likert ба SUS-тэй дүйцүүлсэн тайлбарыг qualitative саналтай хамт уншихыг зөвлөсөн \\
\hline
Comparison validity & Бусад model/system-тэй харьцуулалт ижил нөхцөлд хийгдээгүй & 6.6-д head-to-head benchmark биш гэдгийг ил тод заасан \\
\hline
\end{longtable}

Эдгээр эрсдэлийг тодорхой тайлбарласнаар тайлагнасан өндөр үзүүлэлтүүдийн үнэ цэнэ буурахгүй. Харин эсрэгээр, үр дүнг ямар хүрээнд итгэлтэй уншиж болох, аль хэсгийг цаашид илүү хатуу баталгаажуулах шаардлагатайг тодруулдаг.

\section{Бүлгийн дүгнэлт}

Туршилтын үр дүнгээс дараах дүгнэлтүүдийг хийж болно:

\begin{enumerate}
\item \textbf{Өндөр нарийвчлал:} Сургасан YOLOv8n загвар нь mAP50 хэмжүүрээр 92.3\% нарийвчлал үзүүлсэн нь жижиглэн худалдааны орчинд хэрэглэхэд хангалттай түвшин юм. Гэвч олон random seed бүхий давталттай туршилт хийгдээгүй тул энэхүү дүнг нэг үндсэн сургалтын туршилтын үр дүн гэж тайлбарлах нь зохистой.
\item \textbf{Бодит хугацааны боловсруулалт:} Backend inference pipeline 245ms дундаж хугацаатай байсан нь real-time шаардлагыг хангахуйц түвшин үзүүлсэн.
\item \textbf{Аудитын үр дүнтэй байдал:} Автомат аудитын систем 90\% нарийвчлалтай ажилласан боловч \texttt{FAIL} ангилал detector-ийн False Negative алдаанд илүү мэдрэмтгий байсан. Иймээс audit engine-ийн чанарыг detector-оос тусгаарлан бус, цуваа системийн гүйцэтгэл гэж үнэлэх нь зөв.
\item \textbf{Хэрэглэгчийн сэтгэл ханамж:} 5 онооноос 4.34 дундаж үнэлгээ авсан нь системийн хэрэглээний чадамж өндөр байгааг харуулж байна. Үүний зэрэгцээ status тайлбар, discrepancy breakdown-ыг илүү ойлгомжтой болгох хэрэгцээ qualitative ажиглалтаар илэрсэн.
\end{enumerate}
